\documentclass[instructions.tex]{subfiles}
\begin{document}
\chapter{Ce que le Sauveur a souffert dans son sacré corps pour expier le péché.}


Jésus-Christ a souffert dans son sacré corps plus que tous les martyrs et tous les hommes ensemble n’ont jamais souffert. O justice de Dieu ! que vous êtes impénétrable ! Pourquoi le Père éternel a-t-il donc permis que des hommes traitassent son Fils avec tant de cruauté et de fureur ? Nous le comprendrions, si nous connoissions ce que c’est qu’un Dieu offensé.

L’histoire nous apprend qu’un roi étant attaqué et assiégé dans sa ville capitale par ses sujets révoltés, son fils alla se jeter au pied du trône, et demanda à son père le pardon pour ces misérables. Le père, voyant que son fils s’offroit en caution, et prenoit la défense de ces rebelles, le conduisit sur les murs de la ville, et l’exposa à des tourmens affreux. Il permit que les démons, les puissances des ténèbres s’emparassent de l’esprit des bourreaux pour leur ôter tout sentiment d’humanité envers ce cher Fils\footnote{Ego autem sum vermis et non homo, opprobrium hominum et abjectio plebis. Ps. 21.}. Il permit que son corps adorable fût meurtri, souffleté, déchiré, chargé de plaies ; qu’il fût couronné de cruelles épines et défiguré ; qu’il fût brisé, comme moulu sous les coups \footnote{Haec est hora vestra et potestas tenebrarum. Luc. 22.} ; et qu’enfin, après avoir répandu tout son sang, on lui disloquât les os, qu’on le clouât comme un scélérat, et qu’il expirât sur une croix, comme un agneau entre deux criminels. Pourquoi tant de rigueur envers ce Fils innocent ? C’est parce qu’il s’est fait notre caution, en s’offrant à son Père pour l’expiation de nos péchés. Oh ! si vous savez peser le prix de nos péchés, quels sentimens n’en aurez-vous pas sur le Calvaire !

Apprenez-nous donc, dit saint Pierre, de la pensée des souffrances de Jésus-Christ. Vous y apprendrez trois choses : 1. quelle est la grandeur de Dieu, puisqu’il n’a pu être justement honoré que par les anéantissements de son Fils ; quelle est la sainteté de Dieu et sa haine pour le péché, puisqu’il avoit comme abandonné son Fils, parce qu’il voyoit en lui la figure du péché ; quelle est la justice de Dieu, qui n’a pu être apaisée que par les humiliations de son Fils ; quelle est la miséricorde de Dieu, qui a livré ce Fils adorable à tant de rigueurs, pour vous racheter de l’enfer; et quelle a été la charité du Fils, qui s’est livré aux supplices et à la mort pour votre amour, et avec autant d’amour que si vous aviez été seul au monde.

2. Vous y apprendrez ce que vous êtes ; quelle est la dignité de votre âme, ce qu’elle a coûté et ce qu’elle vaut, puisqu’il a fallu qu’un Dieu suât le sang, mourût dans les tourmens pour la réparer, pour la délivrer et la sauver.

3. Vous apprendrez quelle est la malignité du péché, puisqu’il a causé la mort du Fils de Dieu ; quelle est son énormité, puisque, sans Jésus-Christ, tous les hommes n’auroient jamais pu effacer un seul péché mortel ; quels châtimens il mérite ; car s’il a fallu les mérites infinis d’un Homme-Dieu pour l’expier, il faut dire aussi qu’une peine infinie, un enfer éternel n’est pas trop rigoureux pour le punir. Vous apprendrez enfin ce que vous devez faire pour le punir en ce monde, en disant : Si le Saint des saints a souffert pour des péchés qu’il n’avoit pas commis, à quoi doit s’attendre un scélérat comme moi, coupable de péchés sans nombre\footnote{Si in viridi ligno haec fiunt, in arido quid fiet ? Luc. 23.} ?

C’est donc avec raison que les saints Pères ont dit que Jésus-Christ crucifié est un grand livre, où les savans et les ignorans doivent apprendre leur religion. C’est dans ce livre divin que les saints ont appris à connoître Dieu, et ce qu’ils doivent à Dieu. Le crucifix est toujours devant nos yeux : bien loin d’approfondir les leçons qu’il nous donne, nous imitons les enfans qui manient un livre qu’ils ne comprennent pas ou qu’ils ne lisent point. Cependant quiconque ne sait pas ce livre de la croix, peut dire qu’il ne sait rien. Saint Paul se glorifioit de ne savoir que Jésus crucifié. Ah ! si les larmes, le sang et la mort d’un Dieu ne sont pas capables de nous ouvrir les yeux sur nos devoirs et sur nos péchés, que nous sommes aveugles et que nous sommes misérables !

\section{Résolutions}

\primo Non, rien ne sera plus capable de me faire balancer dans la sérieuse entreprise de ma conversion.  

\secundo Je suis prêt, aidé de la grâce divine, à tout souffrir, plutôt que de retomber dans un seul péché mortel. 

O Jésus ! je vous en conjure par les souffrances que vous avez endurées à cause de mes transgressions, aidez-moi à concevoir une vive horreur du péché, et à faire pénitence de ceux que j’ai eu le malheur de commettre.

\end{document}
